\begin{lemma}{Existenz (unabh\"angiger) F--verteilter Zufallsgrößen}
             {ExistenzFVerteilterZufallsgroessen}
Seien $F: \R \rightarrow [0,1]$ eine Verteilungsfunktion und $n \in \N$.
Dann gilt:
\begin{enumerate}[(1)]
	\item \label{ExistenzFVerteilterZufallsgroessen1}
	Es gibt einen Wahrscheinlichkeitsraum $(\Omega_1, \mfc, \Prim_1)$
	und $\xi \in \mcm(\Omega_1, \mfc, \R)$, sodass $\xi$ $F$--verteilt
	 ist.\\
	Insbesondere ist $\xi$ sicher endlich.
	\item \label{ExistenzFVerteilterZufallsgroessen2}
	Es gibt einen Wahrscheinlichkeitsraum $(\Omega_2, \mfd, \Prim_2)$ und
	$\eta \in \mcm_n(\Omega_2, \mfd, \R^n)$ derart, dass $\eta$ $F$--verteilt  
	und unabhängig ist (im Sinne von~\ref{UnabhaengigIdentischVerteilt}).\\
	Insbesondere ist $\eta_i$ sicher endlich für alle $i \in \haken n$.
\end{enumerate}
\end{lemma}
\begin{proof}
Als Verteilungsfunktion ist 
$F$ monoton wachsend, rechtsseitig stetig und es gilt
\[\Limes t {-\infty}F(t) = 0\text{, sowie }\Limes t \infty F(t) = 1\text.\]
(1) Aus~\cite{Irle}, dort Beweis von 6.6, geht hervor, dass gilt:
\[G:\!\  ]0,1[ \rightarrow \R, p \mapsto \inf\menge{x \in \R}{F(x) \ge p}\]
ist wohldefiniert und genügt der Äquivalenz
\begin{align}
\forall\,t\in\R\,\forall\,s\in\!\ ]0,1[:
G(s)\le t\iff s\le F(t)\text.\label{efvzeq:1}
\end{align}
Seien $\Omega_1\coloneqq \!\ ]0,1[$, $\mfc \coloneqq  \Pot(]0,1[)\: \cap\: \mfb$ und
$\Prim_1\coloneqq \restrict{\lambda}{\mfc}$, wobei $\lambda$ das eindimensionale Lebesguemaß
bezeichne.\\
Sei $t \in \R$.\\
Dann ist
\[\meng{G \le t} = \menge{x \in\!\  ]0,1[}{G(x) \le t}
\overset{\eqref{efvzeq:1}}= 
\menge{x \in\!\  ]0,1[}{x \le F(t)} = \!\ ]0, F(t)] \in \mfc\]
und daher $\Prim_1(G \le t) = \lambda (]0, F(t)]) = F(t)$, sodass
$G$ $F$--verteilt ist.\\
Damit wurde ein Wahrscheinlichkeitsraum und eine Abbildung wie behauptet
gefunden.\absatz
(2) Seien $\Omega_2\coloneqq \!\ ]0,1[^n$, $\mfd\coloneqq \Pot\left(]0,1[^n\right) \cap \mfb_n$
und $\Prim_2\coloneqq  \restrict{\lambda_n}{\mfd}$, wobei $\lambda_n$ das $n$--dimensionale
Lebesguemaß bezeichne.\\
Nach (1) existiert ein $F$--verteiltes $\xi:\!\ ]0,1[\rightarrow\R$.\\
Sei $\eta: \Omega_2 \rightarrow \R^n, \omega \mapsto \pfeil \xi n (\omega)$.\\
Offenbar gilt dann für alle $j \in \haken n$ und alle $\omega \in \Omega_2$:
$\eta_j(\omega) = \xi(\omega_j)$.\\
Wir können o. B. d. A. annehmen, dass $n > 1$ ist, denn für $n = 1$ sagt (2)
dasselbe aus, wie (1).\\
Seien $i \in \haken n$ und $t \in \R$.\\
Falls $1<i<n$, so gilt:
\begin{align*}
\meng{\eta_i \le t} & = \menge{\omega \in \Omega_2}{\eta_i(\omega) \le t}
 = \menge{\omega \in \Omega_2}{\xi(\omega_i) \le t}\\
 &=\!\  ]0,1[^{i-1}\times\menge{x \in\!\  ]0, 1[}{\xi(x)\le t}\times\!\ 
 ]0,1[^{n-i}
 \in \mfd
\end{align*}
falls $i = 1$:
\begin{align*}
\meng{\eta_i \le t} & = \menge{\omega \in \Omega_2}{\eta_1(\omega) \le t}
 = \menge{\omega \in \Omega_2}{\xi(\omega_1) \le t}\\
 &= \menge{x \in\!\  ]0,1[}{\xi(x) \le t}\times\!\  ]0,1[^{n-1} \in \mfd
\end{align*}
und falls $i = n$:
\begin{align*}
\meng{\eta_i \le t} & = \menge{\omega \in \Omega_2}{\eta_n(\omega) \le t}
 = \menge{\omega \in \Omega_2}{\xi(\omega_n) \le t}\\
 &=\!\  ]0,1[^{n-1} \times \menge{x \in \!\ ]0,1[}{\xi(x) \le t}  \in \mfd
\end{align*}
Außerdem gilt
\begin{align}
\Prim_2(\eta_i \le t) & = \lambda_n(\eta_i \le t) \notag\\
&= \begin{cases} 
\lambda_n\bigg(
]0,1[^{i-1}\times\menge{x\in\!\ ]0,1[}{\xi(x)\le t}\times\ \!]0,1[^{n-i} \bigg)
 & : \, 1 < i < n\notag\\
\lambda_n\bigg(
\menge{x \in\!\  ]0,1[}{\xi(x) \le t}\times\!\ ]0,1[^{n-1} \bigg) & : \,i = 1
\notag\\
\lambda_n\bigg(]0,1[^{n-1}\times\menge{x\in\!\  ]0,1[}{\xi(x) \le t} \bigg)&:\,
i=n\end{cases}\notag\\
&= \lambda\left(\menge{x \in\!\  ]0,1[}{\xi(x) \le t}\right) \cdot
 \lambda_{n-1}\left(]0,1[^{n-1} \right)\notag\\
 &= \lambda\left(\menge{x \in\!\  ]0,1[}{\xi(x) \le t}\right)
  \overset{(1)}= F(t) \label{efvzeq:2}
\end{align}
Damit ist $\eta$ $F$--verteilt.\\
Sei nun $s \in \R^n$. Dann gilt:
\begin{align*}
& \Prim_2\left(\bigcap_{k=1}^n \meng{\eta_k \le s_k}\right)
= \Prim_2\left(\forall\, k \in \haken n: \eta_k \le s_k\right)\\
 =&\, \lambda_n\left( \menge{\omega \in \Omega_2}
  {\forall\, k \in \haken n: \eta_k(\omega) \le s_k} \right)
= \lambda_n\left(\menge{\omega \in \Omega_2}
  {\forall\, k \in \haken n: \xi(\omega_k) \le s_k} \right) \\
= &\,\lambda_n \left(\prod_{k = 1}^n \menge{x \in \!\ ]0,1[}
  {\xi(x) \le s_k}\right)
= \prod_{k=1}^n \lambda\left(\menge{x \in\!\  ]0,1[}
  {\xi(x) \le s_k} \right)\\
\overset{(1)}=&\, \prod_{k=1}^n F(s_k) \overset{\eqref{efvzeq:2}}= 
\prod_{k=1}^n \Prim_2(\eta_k \le s_k)
\end{align*}
Damit ist $\eta$ auch unabhängig und die Behauptung gezeigt.
\end{proof}